\chapter{Clause Structure and Word Order}\label{ch:clause-structure}\label{ch:simple}

\section{Theoretical background and analytical conventions}
\label{sec:th-background}

The syntactic analyses in this grammar adopt a Minimalist
architecture in which clause structure is built by Merge, with the restricted
ordered-edge insertion defined in §\,\ref{sec:pivot-placement}, and interpreted through
hierarchies of lexical and functional projections, and linearised according to
the strongly head-final character of Iridian. Recurrent clause properties are
analysed through a shared structural spine comprising a predicate domain,
inflectional and operator projections, and an articulated left periphery.

Within that architecture, the clause edge comprises an ordered sequence of
projections for topic, focus, ordinary placement, clause typing, evidential
and epistemic operators, and speech-act marking. Topic placement, wh-licensing, polar questions,
complementation, and the order of high particles are analysed through these
positions, with information structure encoded especially in
\synhead{TopP} and \synhead{FocP}.

The morphology discussed in this grammar is interpreted in the same way as
elsewhere in the grammar: broadly within a Distributed Morphology model.
Syntactic structure is built first, and the exponence of heads and feature
bundles is determined at spellout. This matters directly for Iridian because
many contrasts that look at first glance like separate `constructions' are
better treated as alternative realisations of the same structural heads:
Neg\textsuperscript{0} surfaces as \ird{za-}, \ird{staly}, or zero-plus-root
suppletion depending on the host available at PF; converb suffixes realise
feature bundles on non-finite heads; and case morphology spells out configurations established in the syntax.

The central descriptive claim of the syntactic analysis is therefore more
specific than a general appeal to `topic-comment structure'. Iridian clause
structure is organised by a head-final functional sequence in which the
absolutive pivot is strongly privileged for topical interpretation, but topic,
focus, polarity, evidentiality, and clause combining are all represented by
independent structural positions. The split-ergative alignment, the ordering of
particles, the placement of subordinate clauses, the distribution of wh-phrases,
and the complement system are analysed as consequences of that shared
architecture. This section sets out the main information-structural
assumptions and terminological conventions used throughout the grammar; the
clause spine itself is given in Table~\ref{tab:th-clause-spine} in this chapter, and the particle framework is analysed in detail in
§\,\ref{sec:th-particle-framework} in Chapter~\ref{ch:function-words}.


\section{Grammatical architecture}\label{sec:grammatical-architecture}
\label{sec:th-clause-structure}% cross-reference compatibility

At a broad level, the Iridian clause is built from the lexical predicate
outward through a sequence of functional layers. The lowest layers belong to
the lexical predicate and its inner valency structure, followed by viewpoint
aspect (Asp.Align); polarity is composed
within the predication selected by phasal operators. Above that domain stand
comment-domain operators, inflection, interrogation, evidentiality, and
speech-act marking; ordinary placement, marked focus and linked topics occupy
distinct positions above the pronounced particle field. High specifiers precede their complements, while the predominantly
head-final predicate domain places the finite verbal complex last. The
interrogative head \ird{by} is locally head-initial. Lexical aspect concerns
the event description in the inner predicate; viewpoint Asp.Align presents that
event, Phas evaluates its polarised proposition across stages, and Mood,
comment operators and evidentiality contribute separate higher meanings.

This architecture is especially visible in the particle system of
Chapter~\ref{ch:function-words}. Particles nearest the predicate express
phasal meanings such as \ird{vy} \trsl{already} and \ird{so} \trsl{still / not
yet}; above these stand comment-domain particles such as \ird{och}
\trsl{also} and \ird{mó} \trsl{indeed}; above these in turn stands
the interrogative clause-typing particle \ird{by}, which is the overt
exponent of a locally head-initial Pol\textsuperscript{0} head above the
inflectional domain; and farther left within the particle field stand evidential and stance
particles such as \ird{oče}, \ird{hladě}, and \ird{kady}. The specified descriptive order
constrains this working clause spine together with morphology and scope; order
alone does not establish a projection for each particle or license multiple
particles within one zone (§\,\ref{sec:particle-combinations}). The narrowly
licensed lower interpretation of focus-associated \ird{kamo}
requires the additional composition rule specified in
§\,\ref{sec:th-exclusive-phasal}.
Within that analysis, \ird{by} realises a dedicated interrogative
clause-typing head in \synhead{PolP}, below the ordinary placement field
\synhead{OrdP}, whereas comment-domain particles occupy lower positions.
The specified partial order is outer speech-act/evidential material $>$
\ird{by} $>$ \ird{och} $>$ \ird{jas} $>$ \ird{vy}/\ird{so};
it does not license every combination of these particles.

The same architecture also helps explain several otherwise disconnected facts.
The final position of the finite verb, the preverbal placement of complements,
the order of subordinate and matrix clauses, and the clause-final position of
\ird{staly} in zero-copula negation all follow from the same head-final
design. Within that structure, aspect and Voice license the case frame,
while a distinct topic dependency organises the comment about a selected
referent. Table~\ref{tab:th-clause-spine} summarises the functional layers
used throughout the grammatical analysis.

The inventory separates thematic argument positions from high discourse
positions. The direct-case pivot is a constructionally selected participant,
whereas Spec,\synhead{TopP} is an independently licensed A$'$-position.
\synhead{IP} abbreviates the finite inflectional domain defined below; the
label does not assert a person-agreement projection or a tense category absent
from the morphology. Neutral pivot-first order in ordinary finite verbal
clauses is derived by the ordinary placement dependency in \synhead{OrdP};
its ordered edge positions do not identify every displaced constituent as
either a topic or an A-moved subject (§\,\ref{sec:pivot-placement}). Reduced predicates and existential constituents
retain the construction-specific structures described in the syntax chapters.

\begin{table}[ht]
	\footnotesize
	\caption{The Iridian clause spine (innermost to outermost).}\label{tab:th-clause-spine}
	\begin{tabularx}{\textwidth}{l l L}
		\toprule
		{\sc level} & {\sc projection} & {\sc content} \\
		\midrule
        1 (lowest) & \synhead{vP}/\synhead{VP} & Lexical predicate and arguments \\
        2 & \synhead{VoiceP}/\synhead{ApplP} & Inner valency and silent actor introduction \\
        3 & \synhead{CausP} & Causation over the inner predicate \\
        4 & \synhead{VoiceP}$_{\mathrm{orient}}$ & Orientation over the derived predicate \\
        5 & \synhead{AspP} & Viewpoint and alignment (Asp.Align) \\
        6 & \synhead{NegP} & Indicative clausal negation: verbal \ird{za-}, equative
          \ird{staly}, or existential \ird{zad-} with zero Neg exponence
          (§\,\ref{sec:vm-neg}) \\
        7 & \synhead{PhasP} & \ird{vy}/\ird{so} over polarised predication \\
        8 & Comm-field & Additive, exclusive, affirmative, contrastive \ird{jas} and interim-sequential operators \\
        9 & \synhead{MoodP} & Mood and finiteness \\
        10 & \synhead{EvidP}$_{\mathrm{inner}}$ & Verbal indirect evidential \ird{-al}/\ird{-l} \\
        11 & \synhead{PolP} & Head-initial interrogative clause typing: \ird{by}/$\varnothing$ \\
        12 & \synhead{EvidP}$_{\mathrm{outer}}$ & Evidential and epistemic operators:
          \ird{oče}, \ird{hladě}, \ird{kady}, \ird{či} \\
        13 & \synhead{SpeechActP} & Illocutionary and discourse framing:
          \ird{nadě}, \ird{mlada} \\
        14 & \synhead{OrdP} & Ordered ordinary placement of arguments and modifiers \\
        15 & \synhead{FocP} & Marked focus or a single marked wh-fronting dependency \\
        16 (highest) & \synhead{TopP} & Optional linked aboutness topic; case retained \\
		\bottomrule
	\end{tabularx}
\end{table}

\noindent This is a working hierarchy with optional domains, not a claim that
every layer is present in every clause. In a verbal clause, \term{predicate
domain} descriptively includes \synhead{AspP} and its inner orientation,
causative, valency and lexical structure; it does not add a separate
\synhead{PredP}. The inner actor-introducing head, where licensed, lies above
Appl and below Caus; its zero exponent does not make it identical to the
verbal categoriser or the outer orientation head
(§\,\ref{sec:vm-overview}). Zero-copula \synhead{EquP} lacks verbal Asp.Align
(§\,\ref{sec:copula}). \synhead{IP} abbreviates the inflectional domain,
including Mood and the inner evidential when present. The Comm-field is a
common descriptive region for heterogeneous operators, not an independently
established single head or a uniform semantic type. The comparison of adverb
order with independently motivated functional-head order follows the kind of
methodological comparison developed by Cinque~\cite[pp.~v--vi]{cinque1999}; Iridian's bounded
working architecture does not adopt his stronger universal hierarchy as an
empirical conclusion about this language.

The table records dominance and operator positions; surface order additionally
 depends on head direction and the positions in which particles are realised.
 Upper-peripheral particles occupy specifiers, while head-initial
 Pol\textsuperscript{0} precedes its complement. Ordinary arguments and
 modifiers are pronounced at the ordered edge of \synhead{OrdP}, above those
 particle positions; linked topics and marked focus have separate higher
 positions. Lower chain occurrences preserve argument structure, case and
 the independently licensed interpretations, including negative concord.
 The lower predicate domain remains head-final
 (§\,\ref{sec:th-head-finality}).

Indicative negation occupies the lower polarity domain, whereas negative
non-indicatives realise polarity through Mood and interpret their negative
operator at the AspP edge. The distribution of exponents and the concord
dependency receive their principal treatment in Chapter~\ref{ch:negation};
phasal and alternative-sensitive composition is developed in
Chapter~\ref{ch:particles}.

The clause fields are schematised in Figure~\ref{fig:th-clause-fields};
Figure~\ref{fig:th-left-periphery} shows the detailed hierarchical structure of
the left periphery.

\begin{figure}[ht]
	\centering
	\begin{tikzpicture}[font=\footnotesize]
\node[syn band, minimum width=27mm] (top)
 {Topic / focus\\\synhead{TopP} / \synhead{FocP}};
\node[syn band, minimum width=29mm, right=0pt of top] (ord)
 {Ordinary placement\\\synhead{OrdP}};
\node[syn band, minimum width=42mm, right=0pt of ord] (part)
 {Articulated particle field\\Speech-act / evidential / \ird{by}\\Comm-field / Phas};
\node[syn band, minimum width=39mm, right=0pt of part] (pred)
 {Predicate / predicator\\Lower argument occurrences\\Inflected verbal complex};

	\end{tikzpicture}
	\caption{Ordinary finite surface fields; nominal predication has construction-specific orders.}\label{fig:th-clause-fields}
\end{figure}

\begin{figure}[ht]
	\footnotesize
	\caption{Hierarchical structure of the Iridian left periphery (\synhead{TopP}
	to \synhead{IP}).}\label{fig:th-left-periphery}
	\medskip
	\centering
\begin{forest}
 syntax tree
 [\synhead{TopP}
  [{DP\\linked topic (optional)}]
  [{Top$'$}
   [\synhead{FocP}
    [{XP\\marked focus (optional)}]
    [{Foc$'$}
     [\synhead{OrdP}
      [{XP$_1$\ldots XP$_n$\\ordinary ordered edge}]
      [{Ord$'$}
       [\synhead{SpeechActP}
        [{Speech-act particle}]
        [{\synhead{EvidP}$_{\mathrm{outer}}$}
         [{Evidential particle}]
         [\synhead{PolP}
          [{Pol\\\ird{by} / $\varnothing$}]
          [\synhead{IP} [{Comm / Phas / lower copies / predicate}]]
         ]
        ]
       ]
       [{Ord$^0$\\$\varnothing$}]
      ]
     ]
     [{Foc$^0$\\$\varnothing$}]
    ]
   ]
   [{Top$^0$\\$\varnothing$}]
  ]
 ]
\end{forest}
\end{figure}


\section{Head-finality}
\label{sec:th-head-finality}

Iridian is predominantly head-final. Its lexical phrases and lower clauses are
internally verb-final or head-final: arguments precede their predicates,
modifiers precede their heads, adpositions follow their nominal complements (as
postpositions, cf.\ §\,\ref{sec:nm-postpositions}), and subordinate clauses precede
the matrix predicate. This property, which is statable as a general principle of
phrase structure, has several major
consequences that recur throughout the syntactic analysis.

For the clause in isolation, head-finality means that the finite verb is always
clause-final in standard Iridian: it is the head of the highest verbal
projection and accordingly surfaces after all its specifiers and complements.
For clause-combining, it means that subordinate clauses precede the matrix
clause whose head they are complements of (see §\,\ref{sec:th-subord-types}
of Chapter~\ref{ch:complex}). For the particle system, specifier realisation places each particle before
its lower complement, whereas ordinary placement raises arguments and
modifiers across the articulated particle projections
(§\,\ref{sec:pivot-placement}). Their resulting adjacency does not make
the particles one phonological word.

Head-finality also has consequences for relativisation and complementation.
Relative clauses precede the head noun they modify; finite complement clauses
precede the matrix predicate that embeds them; and purpose-clause converbs
precede the verb they are oriented toward. These properties are not stipulated
separately but follow from the single head-final phrase structure combined with
the clause-spine in Table~\ref{tab:th-clause-spine}. The interrogative head
\ird{by} is a local exception: Pol\textsuperscript{0} precedes its
complement. Ordinary placement puts overt arguments and modifiers before
this head, yielding ordinary constituents $>$ \ird{by} $>$ lower
predication. The other particles retain their independently specified
specifier realisations and silent final heads.

\section{Word order}\label{sec:word-order}

The traditional descriptive shorthand for Iridian word order is \term{SOV}, but
that label should not be mistaken for the structural primitive of the clause. A
more faithful characterisation is \term{pivot-initial, verb-final}, or more
specifically \term{absolutive-first, verb-final}. The finite verb itself is
rigidly clause-final; the ordering of preverbal nominals is governed by
information structure and alignment rather than by a fixed subject-object
template.

In ordinary active nominative-domain clauses, the agent is the absolutive
pivot, so this pattern resembles SOV. In ordinary ergative transitives,
the same template yields a patient-initial order of the type
\term{absolutive $>$ ergative $>$ verb}. Without a distinct high topic, the neutral order is therefore
`pivot before the remaining arguments, with the finite verb last'.

\subsection{Participant selection and finite placement}
\label{sec:pivot-placement}

The direct-case pivot is identified by its case and the construction's
participant-selection rules, independently of its pronounced position.
Orientation, aspect and valency determine which participant is selected;
directive Mood can override ordinary indicative selection
(§§\,\ref{sec:vm-asp}, \ref{sec:aspect-caus}). Selection, case licensing and
placement are distinct dependencies. A relation established by Agree need not
involve displacement, and an exponent inserted at PF cannot retrospectively
choose a nominal for movement in the syntax.

Ordinary finite arguments and modifiers are pronounced at a clause-local
 placement edge, represented here as \synhead{OrdP}, above the particle
 projections. Its silent head bears an ordered set of placement requirements:
 first the participant selected by the completed Voice--Asp.Align--Mood
 configuration, then the remaining eligible constituents in their lower
 structural order. Each eligible chain bears a placement feature distinct
 from its Case and agreement features, so Case licensing does not discharge
 its placement requirement. The initial search restricts that feature to the
 chain identified by participant selection; the residual searches target
 the remaining chains. Ord discharges the feature at the final ordinary
 edge, while intermediate phase edges only register the chain's accessibility.
 A closer non-selected argument cannot intervene in the first search because
 it does not satisfy that search's selection description. Subsequent searches
 target the closest eligible argument or ordinary modifier that has not yet
 been placed. Operator particles, predicate heads and material internal to
 a nominal or selected clause are ineligible goals. Each later goal tucks
 in below the already established edge occupants. Thus the sequence of
 movements preserves the residual order rather than reversing it through
 repeated outermost adjunction. Ord is a placement head, without a topic,
 focus, case-valuing or independently scope-taking feature; its movement is
 not uniformly classified as A-movement.

The ordinary edge contains hierarchically ordered specifiers, rather than
 treating all its occupants as equidistant. While the projection of a placement
 head remains open, successive movements satisfying its placement requirement
 insert below earlier edge occupants; a higher head cannot subsequently reopen
 that completed projection. This restricted tucking-in has a formal precedent
 in Richards's treatment of multiple movement \cite[§2]{richards2002}.
 It relaxes root-only structure extension: Chomsky explicitly recognises
 tucking-in as a potential principled departure from the No-Tampering
 Condition \cite[p.~141]{chomsky2008}. The Iridian ordered-edge rule is
 adopted with that qualification, rather than derived from Merge alone.
 Case-independent placement activity is an assumption of the Iridian
 analysis: the stronger inactivity condition on completed Case/Agree chains
 discussed by Chomsky \cite[pp.~150--151]{chomsky2008} is not imposed on
 this placement dependency.

The lower thematic structure is preserved: the lexical patient merges with
 V, the actor in its Voice domain, and applied participants and causees in
 their independently licensed positions. The analysis assigns an ordinary
 placement edge requirement to each actor- or causer-introducing head treated
 as a phase within one finite clause. Before transfer, eligible ordinary
 constituents must be registered at its accessible edge: existing eligible
 specifiers remain accessible, and lower goals move there nearest first,
 tucking beneath the earlier edge occupants. Each search targets the highest
 accessible occurrence of a chain, and a chain already registered at that
 particular edge is not selected again through one of its lower copies.
 This non-pivot-selective preparation preserves the original residual order; it precedes the completed
 orientation--Asp.Align--Mood configuration on which Ord's first search
 depends. Only Ord gives the selected participant priority over that residual
 order. A non-phasal head adds no escape step, and the analysis does not infer
 the phasehood of every argument-building layer from surface order
 (§\,\ref{sec:vm-pivot-derivations}). Neither an intermediate edge occurrence
 nor final placement removes the lower actor occurrence required for binding. Finite CP and nominal boundaries delimit separate searches:
 placement of a selected clause as a whole does not extract its arguments.
 A selected finite \ird{po}/\ird{poby} CP moves as a whole to a
 construction-specific external edge of \synhead{OrdP}, preceding its
 pivot and residual occupants. A separate requirement for selected finite complements
 applies after ordinary placement while OrdP remains open;
 it licenses an outer edge rather than another tucked-in ordinary occupant.
 The complement completes its own local structure independently, and its
 internal arguments are not matrix placement goals. Reduced
 purposive complements retain their position relative to the pivot (§\,\ref{sec:vm-po-comp}). Ordinary
 placement orders their whole constituents within those licensed positions.

The following derivations make the ordering requirements explicit. Bracketed
 role letters designate whole constituents, and $\Pi$ designates the licensed
 particle sequence; they are structural schemata, not additional lexical
 examples. The arrows show successive pronounced edge occupants, with lower
 copies left in the thematic configuration throughout.

\begin{description}[leftmargin=2em]
\item[\upshape Nominative-domain transitive.]
 From $[A\ [T\ V]]$, the selected actor moves first; the residual search
 then places $T$ beneath it: $A \longrightarrow A,T\longrightarrow
 A\ T\ \Pi\ V$. For \ird{Marek travě pjaštíme}, this preserves actor-direct
 and patient-indirect case, with neither participant assigned a topic feature.
\item[\upshape Ergative-domain transitive.]
 From the same thematic configuration, the selected patient $T$ moves
 across the non-selected actor, followed by residual $A$:
 $T\longrightarrow T,A\longrightarrow T\ A\ \Pi\ V$.
 Thus \ird{Tak tóm Markuto koupnik} retains patient-direct and
 actor-transitive case. The crossing belongs to the selective first search;
 the residual search cannot independently reverse $T$ and $A$.
\item[\upshape Linked actor topic.]
 After ordinary $T,A$ placement in an ergative transitive, the actor's
 independently licensed aboutness dependency moves it to Spec,\synhead{TopP}:
 $T,A\longrightarrow A_i\ [T\ t_i\ \Pi\ V]$.
 This derives \ird{Markuto tak tóm koupnik} without changing the lower
 actor--patient relation or either case (§\,\ref{sec:tc-marked-topics}).
\item[\upshape Additional participant.]
 A simple applicative selects $B$ from the lower $A>B>T$ configuration:
 $B\longrightarrow B,A\longrightarrow B,A,T\longrightarrow
 B\ A\ T\ \Pi\ V$. The existing \ird{Byl Marcí travě pjaštébime}
 has beneficiary-direct, actor-indirect and theme-indirect case. An ergative
 transitive causative instead selects $T$ from $C>E>T$:
 $T\longrightarrow T,C\longrightarrow T,C,E\longrightarrow
 T\ C\ E\ \Pi\ V$, as in \ird{Travó Marku bylam pjášt-r-an-ik}
 (§\,\ref{sec:aspect-caus}). Causee $E$ remains indirect, and the movement
 must pass through the accessible edges of its non-finite Voice domain.
\item[\upshape Manner modification.]
 With lower $[A\ [M\ [T\ V]]]$ and actor selection, the sequence is
 $A\longrightarrow A,M\longrightarrow A,M,T\longrightarrow
 A\ M\ T\ \Pi\ V$. An omitted actor leaves $M\ T\ \Pi\ V$,
 preserving \ird{Skorně travě pjaštíme}. In an ergative clause, selected
 $T$ moves first, after which residual $A$ and $M$ retain their order;
 pivot priority is distinct from the manner--object relation in the
 nominative-domain example.
\end{description}

Only the highest ordinary occurrence is pronounced unless the constituent
 continues to a separately licensed topic or focus position, in which case
 that higher occurrence is pronounced. Lower occurrences remain available
 for the interpretation and licensing specified below. Null arguments supply
 no overt edge occupant. Neither topicless presentation nor quantificational
 reference suppresses ordinary placement. Participant nominals, reduced
 relatives and converbs retain their own selection and placement systems;
 an AspP-selecting construction does not acquire finite Ord, Mood, Pol or
 stance structure through sharing the same participant-selection rules.

\subsection{Scope connectivity}
\label{sec:pivot-scope-connectivity}

A high pronounced argument retains the case and interpretation licensed
 through its lower chain. In \ird{Deno že zastráda}, the lower occurrence of
 the negative indefinite remains in the domain of the single clausal negation
 (§\,\ref{sec:nm-negative-concord}); ordinary placement does not turn it into
 a referential topic. The actor's lower occurrence likewise remains available
 for the established reflexive dependency involving \ird{svem}
 (§\,\ref{sec:nm-pron-refl}). A chain counts as one argument for case:
 neither an additional silent occurrence nor a higher pronounced occurrence
 adds a case competitor, and case is not recalculated from surface order.
 The exclusion of same-chain occurrences from dependent-case competition has
 a precedent in the framework discussed by Baker and Vinokurova
 \cite[§1]{bakervinokurova2010}; it does not make the timing of Case assignment
 invariant across languages or constructions.
 Wh-variables are interpreted at their lower occurrences under Pol, while
 additive and exclusive association accesses the lower argument variable.
 Quantified interpretations retain their independently specified readings;
 ordinary placement alone grants no scope reversal or new binding relation.
 These are interpretation conditions on Iridian ordinary-placement chains,
 rather than consequences of displacement in general. Fox's distinction
 between chain pronunciation and selective interpretation supplies a
 reconstruction model, while his separate treatments of A- and A$'$-chains
 preclude inferring a uniform reconstruction mechanism from high pronunciation
 alone \cite[pp.~157--159, 192--193]{fox1999}. Lower interpretation is retained
 where the particular Iridian dependency requires it, without obligatory
 reconstruction in every respect.

The use of modifier distribution to constrain functional height has a
comparative model in Cinque \cite{cinque1999}, without requiring a dedicated
projection for every Iridian lexical adverb. The position of a pivot relative
to low manner modification, frequency
quantification, phasal particles and comment-domain particles supplies separate
constraints on its structural height. The two bounded exclusive/phasal readings
in §\,\ref{sec:th-exclusive-phasal} already show why overt precedence cannot
serve as a general rule for semantic scope. Those readings and negative-concord licensing constrain chain interpretation
independently of the finite placement edge.

\subsection{Lexical manner and frequency adverbs in the preverbal field}
\label{sec:lexical-adverb-position}\label{sec:lex-adv}

Lexical manner adverbs originate as low adjuncts to \synhead{vP} or the
 outer \synhead{VP}, while frequency modifiers may attach higher within the
 event description. Their ordinary pronounced occurrences occupy the finite
 placement edge with the arguments, preceding the entire particle field.
 Lower copies preserve their event-modifying interpretations; they acquire
 neither topic status nor stance-operator status through displacement.
 In an active nominative-domain transitive, the successive placement steps in
 §\,\ref{sec:pivot-placement} preserve pivot $>$ manner adverb $>$ object,
 with particles following all three.

\noindent With its pivot unexpressed and no overt particles, an intransitive
clause accordingly has the neutral surface order \term{Adverb $>$ Verb}:

\pex
\a
\begingl
\gla Skorně zemácime.//
\glb quickly run-\Prog{}//
\glft \trsl{I am running quickly.}//
\endgl
\xe

\noindent In an active nominative-domain transitive with an unexpressed
actor and no overt particles, the non-pivot object follows the manner adverb:
\term{Adverb $>$ Object $>$ Verb}. The object retains its lower complement
occurrence for interpretation:

\pex
\a
\begingl
\gla Skorně travě pjaštíme.//
\glb quickly bread.\Ind{} eat-\Prog{}//
\glft \trsl{I am eating bread quickly.}//
\endgl
\xe

An overt actor makes the contribution of pivot selection clearer.
The schematic repair examples below keep the technician, pump and manner
constant while changing the aspectual domain:

\pex\label{ex:order-slot-manner-alignment}
\a
\begingl
\gla \textnormal{[[technician]]}\textsubscript{\Dir{}} skorně \textnormal{[[pump]]}\textsubscript{\Ind{}} \textnormal{[[repair]]}\textsubscript{\Prog{}}.//
\glb technician.\Dir{} quickly pump.\Ind{} repair-\Prog{}//
\glft \trsl{A technician is repairing a pump quickly.}
\lingcontext{work is in progress during an interruption to the water supply}//
\endgl
\a
\begingl
\gla \textnormal{[[pump]]}\textsubscript{\Dir{}} \textnormal{[[technician]]}\textsubscript{\Trs{}} skorně \textnormal{[[repair]]}\textsubscript{\Trs{}.\Pf{}}.//
\glb pump.\Dir{} technician.\Trs{} quickly repair-\Trs{}-\Pf{}//
\glft \trsl{A technician repaired a pump quickly.}
\lingcontext{the repair is presented as one completed operation}//
\endgl
\xe

The progressive places the direct actor before the manner expression and
indirect patient. In the perfective, direct patient selection gives the pump
priority, while the residual actor--manner order is preserved. Neither
order makes \ird{skorně} a stance particle, and the pump's initial position
in (b) does not independently establish an aboutness dependency.

\noindent Frequency adverbs such as \ird{ohly} \trsl{sometimes},
\ird{roca} \trsl{often}, \ird{člíně} \trsl{infrequently}, and \ird{sabo}
\trsl{never} belong to the broad preverbal lexical-adverb field. Their event-quantificational
interpretation differs from low manner modification; membership of that field
does not require attachment at the same structural height. Temporal frames
such as \ird{bych} \trsl{yesterday} locate the event in a time interval and
provide a third, distinct comparison. Even so, they are not members of the rigid particle series. More
leftward placement of lexical adverbs is possible under discourse pressure, but
such ordering belongs to marked information structure and should not be taken
as the neutral pattern. Lexical-adverb placement alone does not determine negative licensing
(§\,\ref{sec:nm-negative-concord}).

Korean case-marked adverbials make the difference between an argument's
case and an adjunct's case especially concrete. In Jou's example, both
the duration expression and the book bear accusative morphology, although
only the book is the selected internal argument:

\pex\label{ex:order-comparative-duration}
\a
\begingl
\glpreamble \textit{Korean} \cite[example~1a, p.~2]{jou2026}//
\gla Mina-ka han sikan-ul chayk-ul ilk-ess-ta.//
\glb Mina-\textsc{nom} one hour-\textsc{acc} book-\textsc{acc} read-\textsc{pst}-\textsc{decl}//
\glft \trsl{Mina read a book for one hour.}//
\endgl
\a
\begingl
\glpreamble \textit{Iridian; schematic}//
\gla \textnormal{[[Mina]]}\textsubscript{\Dir{}} tomem kářime.//
\glb Mina.\Dir{} book.\Ind{} read-\Prog{}//
\glft \trsl{Mina was reading a book.}
\lingcontext{a point within the same one-hour reading session; duration is supplied by the context}//
\endgl
\xe

The Iridian schema represents the reading event with an ordinary direct
actor and indirect patient; its contextual hour supplies no pronounced
durative nominal and consequently establishes no case rule for one.
Jou's account of the Korean durative and multiplicative patterns correlates
adverbial case with subject position, using negative-concord intervention
and subject--event scope as independent diagnostics
\cite[§§4--6]{jou2026}. In Iridian, ordinary placement preserves lower
argument-case licensing, as the repair pair in
(\ref{ex:order-slot-manner-alignment}) illustrates; the shared Korean
accusative marking belongs to a durative construction, whereas
\ird{skorně} supplies a derived manner expression.

Manner adverbs in \ird{-ě} form an open derived class
(§\,\ref{sec:stative-adv}), whereas the frequency items constitute a lexical
subclass independent of that derivation. Neither belongs to the closed
particle series. Negative-concord temporal \ird{žemě} \trsl{never}
requires local negation (§\,\ref{sec:nm-negative-concord}); the distribution
of a lexical frequency modifier must be distinguished from that licensing
dependency.


\section{Verbal complements and head finality: preverbal position}
\label{sec:vm-comp-headfinal}

Internal nominal arguments and selected clauses originate in the complement domain
of the head-final predicate. Their lower positions determine selection and
interpretation, while the ordinary finite placement dependency determines
pronunciation before the articulated particle field
(§\,\ref{sec:pivot-placement}).

\subsection{The base position of the object}
\label{sec:vm-obj-base}

In a head-final language, the complement of a head $H$ is to the
left of $H$ in the linear string. For the lexical predicate V, its
internal argument (patient \synhead{DP}) is the complement of V and therefore
precedes V in the lower structure:

\begin{center}
\term{[\synhead{VP} [\synhead{DP} patient] V]}
\end{center}

\noindent This is the \emph{base-generated} complement position. Its
preverbal orientation follows from head-finality; its ordinary pronunciation
before particles additionally involves finite ordinary placement. In the Minimalist
framework used here, head-finality is the linearisation convention assigning
$<$Complement, Head$>$ order in the lower predicate. This directly represented
complement order is an architectural assumption; the further position before
particles follows from ordinary placement, rather than from head-finality
alone.

The patient \synhead{DP} is base-generated as the complement of V, before
 the verb. Its direct or indirect case is licensed within the aspectual and
 verbal domain; case assignment does not require movement to
 Spec,\synhead{TopP}. The patient's origin below the actor does not by itself
 derive its neutral position before that actor. Ordinary placement first
 attracts the direct-case pivot and then the eligible residual constituents
 to ordered positions below it in \synhead{OrdP}; its lower occurrence
 preserves the licensed case and thematic relation. This dependency is
 distinct from linked topicalisation, so a pre-particle pivot can lack an
 aboutness interpretation (§\,\ref{sec:tc-quantificational-pivots}).

When an argument is independently selected as a linked topic, an
 A$'$-dependency connects Spec,\synhead{TopP} with that lower position.
 Schematically, with intervening projections suppressed:
\[
 [_{\mathrm{TopP}}\ \mathrm{DP}_{i}\
   [\mathrm{lower\ clause}\ \ldots t_i\ldots V]\ \mathrm{Top}^{0}].
\]
The \synhead{DP} retains its case, so a perfective transitive can have either
 a direct-case patient topic or a transitive-case agent topic. Head-finality
 places complements before V; the topic dependency, when present, places its
 referent above the lower clause. These are distinct sources of preverbal
 order.

\section{The preverbal zone: a structural hierarchy}
\label{sec:vm-preverbal-zone}

The ordinary surface fields are represented in Figure~\ref{fig:th-clause-fields}.
Selected finite complements and reduced purposive complements retain their
construction-specific positions (§§\,\ref{sec:vm-po-comp},
\ref{sec:vm-purp-comp}), while their internal arguments belong to their
own selection and interpretation domains.

\section{Verbal exponents within the clause spine}
\label{sec:vm-summary}

Table~\ref{tab:vm-full-structure} gives the hierarchical structure
of the Iridian finite verbal clause, from the highest projection
(\synhead{TopP}) to the lowest (\synhead{VP}/Root), with the morphological
exponents of each relevant layer. This working architecture includes optional
domains; it does not posit every projection in every clause. The descriptive
Comm-field groups semantically heterogeneous particles without committing to
a single head. \synhead{IP} elsewhere abbreviates the inflectional domain
(Mood and, when present, the inner evidential); the predicate domain comprises
\synhead{AspP} and its inner structure, without an extra PredP projection.

\begin{table}[ht]
  \footnotesize
  \caption{The Iridian verbal clause: functional hierarchy and
  morphological exponents.}\label{tab:vm-full-structure}
  \begin{tabularx}{\textwidth}{l l l L}
    \toprule
    {\sc level} & {\sc projection} & {\sc morphological exponent}
    & {\sc notes} \\
    \midrule
    Highest & \synhead{TopP} & [topic \synhead{DP}]
    & Linked aboutness topic; lower case retained \\
    & \synhead{FocP} & [marked focus / wh-phrase]
    & At most one marked wh-fronting dependency \\
    & \synhead{OrdP} & $\varnothing$; ordered edge constituents
    & Ordinary arguments and modifiers before particles \\
    & \synhead{SpeechActP} & Outer framing particles
    & Speech-act contribution within the articulated field \\
    & \synhead{EvidP}$_\text{outer}$ & \ird{oče, hladě, kady, či}
    & Clause-level evidential/epistemic operators \\
    & \synhead{PolP} & \ird{by}/$\varnothing$
    & Locally head-initial interrogative clause typing \\
    \midrule
    & \synhead{EvidP}$_\text{inner}$ & \ird{-al}, postvocalic \ird{-l}
    & Indirect evidential; inside verbal complex \\
    & \synhead{MoodP} & \ird{-ím, -é, -ká, -ku, -ěv, -ir}, $\varnothing$
    & Negative Mood bears $[u\textsc{neg}]$; reduces Asp features \\
    & Comm-field & \ird{och, jas, kamo, mó, dlí}
    & Heterogeneous operators; specified order \ird{och} $>$ \ird{jas} $>$ Phas \\
    & \synhead{PhasP} & \ird{vy, so}
    & Selects positive or negative predication \\
    & \synhead{NegP} & \ird{za-} (prefix)
    & Indicative only; Op$_{\neg}$ in specifier; Neg$^{0}[u\textsc{neg}]$; left-edge realisation \\
    & \synhead{AspP} & \ird{-ek/-ik, -av, -a, -ime, -áž} (aspect + alignment)
    & Obligatory finite head; viewpoint and alignment domain \\
    \midrule
    & \synhead{VoiceP}$_{\mathrm{orient}}$ & \ird{-n-}/\ird{-an-}, $\varnothing$
    & Selects participant of the complete derived predicate \\
    & \synhead{CausP} & \ird{-r-}, \ird{-er-}, \ird{-ěr-}
    & Adds causer above the inner valency domain \\
    & Inner actor layer & $\varnothing$
    & Introduces the lexical actor above inner valency \\
    & \synhead{VoiceP} / \synhead{ApplP} & \ird{-aš-} (antipassive) / \ird{-éb-} (applicative)
    & Distinct heads; language-specific co-occurrence ban \\
    & \synhead{vP}/\synhead{VP} & Root stem
    & Lexical predicate; complement of V = preverbal \\
    Lowest & \synhead{VP} complement & \synhead{NP} / \synhead{CP} (\ird{po}) / \synhead{VoiceP}$_{\text{-it}}$
    & Object, po-clause, or purposive converb; all preverbal \\
    \bottomrule
  \end{tabularx}
\end{table}

\noindent The preverbal position of all verbal complements --- \synhead{NP}
objects, \synhead{CP} complements (\ird{po}-clauses), and purposive converb
complements --- follows uniformly from the head-final parameter applied at the
\synhead{VP} level. The complement originates to the left of V; ordinary finite pronunciation
before the particle field additionally requires the placement dependency,
whose lower chain occurrences preserve selection and case. Selected clauses
retain their construction-specific positions as whole constituents.
Asp, Mood and the inner evidential have suffixal exponents on the verbal
complex, which occupies the right edge of the lower clause. Neg is structurally
higher than Asp but attaches at the left edge of the verbal host, as specified
in §\,\ref{sec:vm-neg}.

%% ----------------------------------------------------------------
%% END: verb-minim.tex
%% ----------------------------------------------------------------
